Training a 3DGS scene starts with a point cloud which we take as the initial positions for 3D Gaussians. After that forward and backpropagation steps similar to those in neural networks \cite{thimm2025ML} alternate, modifying opacity, position, shape and color of each 3D Gaussian in the scene.
Pivotal for training quality is so-called \textit{Adaptive Density Control} which periodically steps in after backpropagation, further modifying 3D Gaussians.

In this chapter we will closely follow Kerbl et al \cite{kerbl3DGaussianSplatting2023} and their GitHub repository \cite{kerblRepo}. The entire training procedure is based on neural rendering concepts introduced by Yifan et al \cite{yifanDifferentiableSurfaceSplatting2019} and ''Pulsar'' by Lassner and Zollhöfer \cite{lassnerPulsarEfficientSpherebased2021} which can be seen as the direct predecessor of the rendering pipeline used for 3DGS. However, there are only very few papers that directly disect the pipeline math (like \cite{huang2DGaussianSplatting2024} or \cite{yuMipSplattingAliasFree3D2024} when improving the algorithm) as most papers use 3DGS for various applications. This is why much content in this chapter originates from in-depth websites like the formidable MrNeRF Github repository \cite{githubMrNerf} which provides further links.

\subsection{Point Cloud Initialization}
Initialization in 3D reconstruction is a crucial step as the positions of the initial points or primitives decide how well the algorithm can reconstruct a scene.

In order for the 3DGS learning algorithm to work, 3DGS needs initial geometry which is commonly a set of initial 3D Gaussians. This initial set is derived from input imagery such as photos or frames of a video. If the camera poses of this input imagery are not given directly, which is usually the case, the camera poses have be derived dynamically from the imagery.
A camera pose is designated by a $4\times 4$ \textit{extrinsic matrix} $W$ \cite{marschnerFundamentalsComputerGraphics2021} consisting of the $3\times 3$ rotational matrix $R$ (or $W_{rot}$) which contains the orientation of a camera in 3D space, a $3\times 1$ translation vector $t$ denoting the position of the camera in worldspace and a padding row to make $W$ square: $$W=\begin{bmatrix} R_{11} & R_{12} & R_{13} & t_x \\ R_{21} & R_{22} & R_{23} & t_y \\ R_{31} & R_{32} & R_{33} & t_z \\ 0 & 0 & 0 & 1 \end{bmatrix}$$
This matrix will become important later during forward and backpropagation steps.

The other component of a camera pose is the \textit{intrinsic matrix} which holds parameters of the camera lens like the focal length.

The defacto standard concept in 3D imaging to derive these two matrices for each input image is called \textit{Structure from Motion} (SfM) \cite{ullmanInterpretationStructureMotion1979} offered through a software called \textit{COLMAP}. SfM solves a large geometric optimization problem which finds common keypoints across all input imagery and uses them to provide highly accurate camera poses and a 3D point cloud of the scene. 

3DGS as originally conceived by Kerbl et al uses these points as the origins for the initial 3D Gaussians which are being initialized as perfect spheres with radii deduced from the three neighboring spheres with an opacity of around $0.1$.
The initial color can be any RGB value.

Consequently, having to derive camera poses for initialization first adds significantly to the entire training process which is why other, faster COLMAP-free methods have been proposed, see chapter \ref{ch:improvements}. Also, the authors mention that using random initialization of all 3D Gaussians yields good quality results even though training times are longer.

Finally, the natural logarithm is applied to all initial attribute values of a 3D Gaussian to be compatible with the initial forward pass.

